%% ch09.tex — Chapter 9: Data Visualization at Scale
%% Part of "Data Engineering" textbook

\chapter{Data Visualization at Scale}
\label{ch:visualization}

\epigraph{\textit{``Above all else, show the data.''}}{Edward R.\ Tufte,
\textit{The Visual Display of Quantitative Information},
Graphics Press, 1983}

\vspace{4pt}

\begin{chapterlearning}
After completing this chapter, you will be able to:
\begin{enumerate}
  \item Explain why Anscombe's Quartet demonstrates that summary statistics
        alone are insufficient for data understanding, and describe the
        role of visualization as a primary data quality check.
  \item Define Tufte's data-ink ratio, compute it for a given chart, and
        identify specific chartjunk elements (heavy grids, 3-D effects,
        gradients) that should be removed to maximise it.
  \item Calculate Tufte's lie factor for a chart with a truncated axis and
        explain the distortion it introduces in the viewer's perception of
        the data effect.
  \item Describe Tufte's small multiples and sparklines and explain why
        each is preferable to animation for comparing trends across many
        subsets.
  \item Apply Shneiderman's three-step visual information seeking mantra
        (overview first; zoom and filter; details on demand) to the design
        of an interactive big data dashboard.
  \item Use Shneiderman's seven data type taxonomy to select the appropriate
        visualization type for a given analytical question (temporal, 2D
        geographic, multi-dimensional, tree, and network data).
  \item Describe three techniques for addressing overplotting in
        million-point scatterplots---alpha transparency, hexbin aggregation,
        and kernel density estimation---and explain what each preserves
        that the overplotted chart loses.
  \item Explain how pre-aggregated data cubes enable sub-second interactive
        query response on billion-row fact tables, and identify the latency
        versus freshness trade-off.
  \item Design a streaming visualization architecture using Kafka, Spark
        Structured Streaming, a time-series database, and Grafana, and
        explain when to use WebSocket push versus periodic poll.
  \item Compare six visualization tools---Tableau, Power BI, D3.js,
        Grafana, Apache Superset, and Kibana---on audience, data source,
        scale, and real-time capability, and select the appropriate tool
        for a given scenario.
\end{enumerate}
\end{chapterlearning}

\bigskip

Chapter~\ref{ch:pipelines} established how data arrives in the platform ---
through Kafka pipelines, data lakes, and stream processing architectures.
This chapter asks what happens at the other end: once data has been
ingested, stored, and aggregated, how should it be communicated to the
humans who must act on it?

The answer is not as obvious as it might appear. A petabyte-scale analytical
result compressed into a single number is not trustworthy; a table of ten
million rows is unreadable; and a chart that presents data incorrectly is
worse than no chart at all. The discipline of data visualization provides
a rigorous body of principles---grounded in perceptual psychology,
information theory, and graphic design---that governs how data should be
encoded visually so that human pattern recognition can operate effectively.

This chapter covers the three intellectual pillars of that discipline:
Edward Tufte's \emph{data-ink} theory of minimal graphical design;
Ben Shneiderman's \emph{task-by-data-type taxonomy} and interaction model;
and Mike Bostock's \emph{D3.js}, which made both principles implementable
in interactive web graphics. It then addresses the three specific engineering
challenges that big data introduces for visualization---overplotting, latency,
and streaming---and maps the tool landscape across which practitioners
must choose.

%%====================================================================
\section{Why Visualization Matters: From Summary to Insight}
\label{sec:viz-why}
%%====================================================================

The case for visualization is not merely aesthetic. It rests on a
demonstrable limitation of summary statistics: different data with
radically different structures can produce identical numerical summaries.

%--------------------------------------------------------------------
\subsection{Anscombe's Quartet}
\label{subsec:anscombe}
%--------------------------------------------------------------------

In 1973, statistician Francis Anscombe constructed four datasets---now
called \term{Anscombe's Quartet}---each with 11 observations, each
sharing the following summary statistics to two decimal places:
mean of $x = 9.00$; mean of $y = 7.50$; variance of $x = 11.00$;
variance of $y = 4.12$; correlation $r = 0.816$; regression line
$\hat{y} = 3.00 + 0.500x$.

Despite being statistically indistinguishable, the four datasets describe
four qualitatively different phenomena:
\begin{itemize}
  \item \textbf{Dataset I} is a linear relationship with normally
        distributed residuals---the ``textbook'' regression case.
  \item \textbf{Dataset II} is a perfect curved (quadratic) relationship;
        a linear fit is systematically wrong at every point.
  \item \textbf{Dataset III} is a near-perfect linear relationship with
        one outlier that drags the fitted slope away from the true slope.
  \item \textbf{Dataset IV} has all $x$-values equal to 8 except one;
        a vertical cluster of points with a single leverage point.
\end{itemize}

Anscombe's point was that no data analyst should trust a regression result
without first plotting the data. The quartet has been the standard opening
of every serious statistics course since 1973 and remains the canonical
argument for prioritising visualization over numerical summary.

\begin{keyinsight}{Visualization as the First Quality Check}
At big data scale, Anscombe's problem becomes more severe, not less.
A billion-row table may contain multiple overlapping sub-populations
with opposing trends that cancel in the aggregate mean. Log-scale
distributions (income, city populations, web traffic volumes) make
arithmetic means uninformative. Seasonal effects can be exactly
cancelled by counter-seasonal confounders, leaving a flat year-over-year
mean that masks both effects. Visualization is therefore not a
presentation step that happens after analysis---it is the first and
most important analysis step.
\end{keyinsight}

%--------------------------------------------------------------------
\subsection{Visualization in the Big Data Context}
\label{subsec:viz-bigdata}
%--------------------------------------------------------------------

Big data introduces challenges that do not arise with small datasets:
\begin{enumerate}
  \item \textbf{Scale}: rendering 10 million points in a standard scatter
        plot produces an uninformative solid blob. Visualization must be
        redesigned for scale.
  \item \textbf{Velocity}: streaming data generates new observations every
        second. Charts must update continuously without requiring the user
        to manually refresh.
  \item \textbf{Variety}: heterogeneous data types---geographic coordinates,
        time series, network graphs, free text---require different chart
        types with different interaction paradigms.
  \item \textbf{Latency}: a business analyst querying a 100-billion-row
        Hive table expects a dashboard filter response in under one second.
        The underlying computation may take minutes.
\end{enumerate}

The remainder of this chapter addresses each challenge in turn, building
from foundational principles to engineering solutions.

%%====================================================================
\section{Tufte's Principles of Analytical Design}
\label{sec:tufte}
%%====================================================================

Edward R.\ Tufte is Professor Emeritus of Statistics, Graphic Design, and
Political Economy at Yale University. His 1983 book \textit{The Visual
Display of Quantitative Information}---rejected by every publisher and
self-published after Tufte mortgaged his house to pay for printing---sold
40{,}000 copies in its first year and has been cited over 20{,}000 times.
It is required reading at Harvard, Stanford, Yale, and MIT in data science
and statistics courses, and its principles are embedded in every modern
dashboard tool.

%--------------------------------------------------------------------
\subsection{The Data-Ink Ratio}
\label{subsec:data-ink}
%--------------------------------------------------------------------

Tufte's central concept is the \term{data-ink ratio}:

\begin{defbox}{Data-Ink Ratio (Tufte, 1983)}
The \textbf{data-ink ratio} of a graphic is the proportion of the total
ink used in the graphic that encodes actual data:
\[
  \text{Data-ink ratio} =
  \frac{\text{ink devoted to data}}{\text{total ink used in graphic}}
\]
\textbf{Principle}: maximise the data-ink ratio. The maximum is 1: every
mark encodes data and no mark is decorative. Erasing a non-data mark
that does not reduce information is called \emph{erasing chartjunk}.
\end{defbox}

\term{Chartjunk} is Tufte's umbrella term for ink that does not encode
data. He identifies three categories:
\begin{itemize}
  \item \textbf{Vibrations}: optical illusions caused by dense cross-
        hatching, competing patterns, or high-contrast decorative elements.
  \item \textbf{Grids}: heavy gridlines that compete with the data.
        Tufte recommends light, thin gridlines or none at all; the data
        should dominate the visual field, not the scaffolding.
  \item \textbf{Ducks}: graphic elements that represent themselves rather
        than data. A bar chart drawn as a stack of coins, a temperature
        chart drawn as a thermometer, or any 3-D perspective added to a
        2-D bar chart are all ``ducks'' --- they encode the same information
        as the underlying simple chart but add ink without adding data.
\end{itemize}

\textbf{Practical application.} When designing or reviewing a chart, ask:
``If I erased this element, would the viewer lose information?'' If the
answer is no, erase it. Specific elements that almost always fail this test:
\begin{itemize}
  \item 3-D effects on bar charts, pie charts, or line charts.
        Three-dimensional perspective distorts quantitative comparison:
        the apparent depth makes a bar appear larger than its height warrants.
  \item Gradient fills. A bar coloured dark-at-bottom and light-at-top
        implies a quantitative dimension that does not exist.
  \item Redundant legends when the chart can be directly labelled.
  \item Background colours, shadows, and rounded decorative borders.
\end{itemize}

%--------------------------------------------------------------------
\subsection{The Lie Factor}
\label{subsec:lie-factor}
%--------------------------------------------------------------------

\begin{defbox}{Lie Factor (Tufte, 1983)}
The \textbf{lie factor} of a graphic measures how much the chart
exaggerates or understates the effect it purports to show:
\[
  \text{Lie factor} =
  \frac{\text{size of effect shown in graphic}}
       {\text{size of effect in data}}
\]
A lie factor of 1 is truthful. A lie factor $>1$ overstates the data
effect; a lie factor $<1$ understates it. Lie factors greater than
1.05 or less than 0.95 are considered deceptive by Tufte.
\end{defbox}

The most common source of a high lie factor is a \emph{truncated axis}:
a bar chart whose $y$-axis starts at 80 rather than 0. If two bars
represent values of 82 and 90, the true ratio is $90/82 = 1.10$
(a 10\% difference). If the $y$-axis starts at 80, the bars appear to
have heights of 2 and 10---a ratio of 5---making the same 10\% difference
look like a $5\times$ difference. The lie factor is $5/1.1 \approx 4.5$.

\begin{caution}{The Truncated Axis in Political and Business Reporting}
Truncated axes are ubiquitous in political advertising, corporate
earnings presentations, and cable news graphics. They are technically
not incorrect---the tick marks do state the true values---but they
exploit the human perceptual system's tendency to judge bar charts by
height rather than by reading the axis labels. In big data dashboards,
the most common violation is a live metrics chart where the $y$-axis
autoscales to the data range: a stable metric that fluctuates between
99.98\% and 99.99\% availability appears to swing dramatically between
\textasciitilde{}0\% and \textasciitilde{}100\% of the chart's visual height. Always verify that
availability and ratio charts start their $y$-axes at zero.
\end{caution}

%--------------------------------------------------------------------
\subsection{Small Multiples and Sparklines}
\label{subsec:small-multiples}
%--------------------------------------------------------------------

Two of Tufte's most practically useful design patterns address the problem
of comparing many subsets simultaneously without animation.

\begin{defbox}{Small Multiples (Tufte, 1983)}
\textbf{Small multiples} (also called \emph{trellis charts} or
\emph{facet charts}) are a grid of charts that repeat the same visual
structure---same axes, same scale, same chart type---across different
subsets of the data (different time periods, geographic regions, product
categories, or experimental conditions). The viewer's eye detects
differences and patterns across the grid instantaneously, leveraging
the perceptual system's ability to perform parallel comparison.
The critical requirement is that \textbf{all panels share the same scale};
if panels have different scales, comparison is invalid.
\end{defbox}

\begin{defbox}{Sparklines (Tufte, 2006)}
A \textbf{sparkline} is a small, word-sized data graphic embedded
directly in text or in a table cell alongside the numeric value it
contextualises. A sparkline adds temporal context to a number without
requiring a separate figure. Grafana's stat panel sparkline, the GitHub
contribution heatmap, and Bloomberg terminal trend indicators are all
direct applications of this concept. The design rule is minimal: no axis
labels, no legend, no gridlines---just the trend shape in a space no
larger than the surrounding text.
\end{defbox}

Small multiples are preferable to animation for the following reason:
animation requires the viewer to \emph{remember} the previous frame to
compare it with the current frame. A small-multiples grid allows
simultaneous comparison of all subsets with no memory burden. A classic
example is the New York Times COVID-19 tracker, which showed 50-state
case charts in a small-multiples grid---allowing readers to compare New
York, Texas, and California at a glance---rather than an animated map
that required readers to hold the April 2020 pattern in working memory
while watching the July 2020 pattern unfold.

%%====================================================================
\section{Shneiderman's Visual Information Seeking Mantra}
\label{sec:shneiderman}
%%====================================================================

Ben Shneiderman is Professor of Computer Science at the University of
Maryland and the inventor of the treemap (1991). His 1996 paper ``The Eyes
Have It: A Task by Data Type Taxonomy for Information Visualizations,''
presented at the IEEE Symposium on Visual Languages, has been cited over
5{,}000 times and provides the standard framework for designing interactive
visualisation systems.

%--------------------------------------------------------------------
\subsection{The Three-Step Mantra}
\label{subsec:mantra}
%--------------------------------------------------------------------

Shneiderman's central contribution is the \term{visual information seeking
mantra}:

\begin{keyinsight}{Shneiderman's Visual Information Seeking Mantra}
\begin{center}
\large\textit{``Overview first, zoom and filter, then details on demand.''}
\end{center}
This three-step model describes the \textbf{universal interaction sequence}
for any well-designed interactive visualisation:
\begin{enumerate}
  \item \textbf{Overview first.} Present the complete dataset or the
        highest-level aggregation as the initial view. The user must
        be able to see the whole picture before selecting a subset.
  \item \textbf{Zoom and filter.} Allow the user to narrow the view to
        a region of interest---either spatially (zoom on a geographic
        map, zoom on a time-axis), or by attribute (filter to a specific
        category, date range, or threshold).
  \item \textbf{Details on demand.} Allow the user to inspect the raw
        record for any visible item by clicking or hovering. The tooltip
        or drill-down panel reveals the full record without switching context.
\end{enumerate}
Every major dashboard tool implements this mantra: Tableau's drill-down,
Grafana's time-range zoom, Google Maps' scroll-to-zoom, and Kibana's
faceted search all instantiate the same three-step pattern.
\end{keyinsight}

%--------------------------------------------------------------------
\subsection{Seven Data Types and Their Visualizations}
\label{subsec:seven-types}
%--------------------------------------------------------------------

Shneiderman identifies seven data types, each with a natural family
of visualization forms and a natural set of interaction tasks.

\begin{table}[htbp]
\centering
\renewcommand{\arraystretch}{1.4}
\begin{tabular}{lp{3.6cm}p{4.2cm}}
  \toprule
  \textbf{Data Type} & \textbf{Typical Analytical Question}
    & \textbf{Recommended Visualizations} \\
  \midrule
  \textbf{1-D (linear)}
    & Which category is largest? How are values distributed?
    & Bar chart, histogram, dot plot \\
  \textbf{2-D (planar/geographic)}
    & How is X distributed across space?
    & Choropleth map (rate), proportional bubble map (absolute count) \\
  \textbf{3-D (volumetric)}
    & What is the internal structure of a volume?
    & Volume rendering, 3-D scatter (scientific visualisation) \\
  \textbf{Temporal}
    & How did X change over time?
    & Line chart, area chart, calendar heatmap \\
  \textbf{Multi-dimensional}
    & Any pattern across $>$3 variables?
    & Scatter plot matrix, parallel coordinates, heatmap with clustering \\
  \textbf{Tree}
    & What is the hierarchical breakdown?
    & Treemap, sunburst diagram, dendogram \\
  \textbf{Network}
    & How are entities connected?
    & Force-directed graph, chord diagram, adjacency matrix \\
  \bottomrule
\end{tabular}
\caption{Shneiderman's seven data types with recommended visualisations.
         Matching the chart type to the data type is the most consequential
         design decision in dashboard construction.}
\label{tab:seven-types}
\end{table}

\textbf{The most abused chart.} The pie chart should only be used when
there are at most four categories and the question is ``what fraction of
the whole does each category represent?'' For any comparison question,
any ranking question, any time-series question, or any analysis with
more than four categories, a horizontal bar chart ordered by value
is always more readable. Human perception judges the length of a bar
far more accurately than the angle or arc length of a pie segment---a
finding confirmed by Cleveland and McGill's perceptual accuracy studies
(1984), which are the empirical foundation of both Tufte's and
Shneiderman's recommendations.

\begin{table}[htbp]
\centering
\renewcommand{\arraystretch}{1.4}
\begin{tabular}{lp{3.2cm}p{4.5cm}p{2.0cm}}
  \toprule
  \textbf{Data Type} & \textbf{Question Type}
    & \textbf{Best Visualization} & \textbf{Avoid} \\
  \midrule
  Temporal    & Trend over time  & Line chart, area chart, calendar heatmap & Pie chart \\
  Categorical & Which is largest? & Horizontal bar chart ($>$4 cats); grouped bar & 3-D pie \\
  Geographic  & Spatial distribution & Choropleth (rate); bubble map (count) & Plain bar \\
  Relational  & Entity connections & Force-directed graph; chord diagram & Table \\
  Distribution & Shape and spread & Histogram; box plot; violin plot & Line chart \\
  Correlation & Does X predict Y? & Scatterplot; scatterplot matrix & Bar chart \\
  Hierarchical & Part-of-whole breakdown & Treemap; sunburst & Nested pie \\
  \bottomrule
\end{tabular}
\caption{Visualization selection by data type and question. The ``Avoid''
         column lists common mistakes that produce low-information charts
         for the given combination of data type and question.}
\label{tab:chart-selection}
\end{table}

%%====================================================================
\section{Big Data Visualization Challenges}
\label{sec:bigdata-viz-challenges}
%%====================================================================

%--------------------------------------------------------------------
\subsection{Challenge 1: Overplotting at Scale}
\label{subsec:overplotting}
%--------------------------------------------------------------------

A standard scatterplot renders one point per data record. At ten million
records, points overlap completely, producing a solid blob that conveys
no more information than ``data exists in this region.'' This is the
\term{overplotting} problem, and it appears whenever visualising a large
dataset with more than $\sim$10{,}000 points in a 2-D projection.

Three solutions exist in increasing sophistication.

\textbf{1. Alpha transparency.} Set each point's opacity to a small value
(e.g., $\alpha = 0.01$). Regions where points overlap appear darker because
the opacities stack multiplicatively. This is simple to implement and
reveals gross density differences, but fails at very high densities: once
$>$100 points overlap, all regions appear equally saturated at $\alpha = 0.01$.

\textbf{2. Hexbin aggregation.} Divide the 2-D plot area into a regular
hexagonal grid (hexagonal cells tile a plane more uniformly than square
cells). Count the number of data points falling in each hexagonal cell.
Colour each cell by its count using a sequential colour scale (e.g.,
light-to-dark blue). The result converts a 10-million-point scatter into
a grid of $\sim$10{,}000 hexagonal cells, each representing a local density
estimate---interactive and readable.

\begin{lstlisting}[style=pyspark, caption={Hexbin scatterplot in Python (matplotlib)}]
import numpy as np
import matplotlib
matplotlib.use("Agg")
import matplotlib.pyplot as plt

rng = np.random.default_rng(42)

# Simulate 5 million GPS coordinates around Dhaka
x = rng.normal(90.41, 0.15, 5_000_000)   # longitude
y = rng.normal(23.81, 0.10, 5_000_000)   # latitude

fig, axes = plt.subplots(1, 2, figsize=(14, 5))

# Left: raw scatter (overplotted)
axes[0].scatter(x, y, s=0.01, alpha=0.003, color="steelblue")
axes[0].set_title("Scatterplot (5M points — overplotted)")
axes[0].set_xlabel("Longitude")
axes[0].set_ylabel("Latitude")

# Right: hexbin (density readable)
hb = axes[1].hexbin(x, y, gridsize=80, cmap="Blues", mincnt=1)
fig.colorbar(hb, ax=axes[1], label="Trip count")
axes[1].set_title("Hexbin (same 5M points — density readable)")
axes[1].set_xlabel("Longitude")
axes[1].set_ylabel("Latitude")

plt.tight_layout()
plt.savefig("dhaka_gps_hexbin.png", dpi=150, bbox_inches="tight")
print("Saved dhaka_gps_hexbin.png")
\end{lstlisting}

\textbf{3. Kernel Density Estimation (KDE).} A \term{kernel density
estimate} places a smooth kernel function (typically Gaussian) on each
data point and sums the kernels to produce a continuous estimate of the
probability density over the 2-D domain. The density surface is then
visualised as filled contours or as a heatmap. KDE reveals the shape,
modes (peaks), and tails of the point distribution more clearly than
hexbin, but is more computationally expensive: $O(N \cdot G)$ where $N$
is the number of points and $G$ is the number of grid cells used to
evaluate the density.

\begin{lstlisting}[style=pyspark, caption={2-D KDE density plot (scipy + matplotlib)}]
from scipy.stats import gaussian_kde
import numpy as np
import matplotlib.pyplot as plt
import matplotlib
matplotlib.use("Agg")

rng = np.random.default_rng(0)
x = rng.normal(90.41, 0.12, 200_000)
y = rng.normal(23.81, 0.08, 200_000)

# KDE on a random subsample (full dataset is too large for naive KDE)
idx    = rng.choice(len(x), size=50_000, replace=False)
kde    = gaussian_kde(np.vstack([x[idx], y[idx]]),
                      bw_method=0.05)

# Evaluate on a grid
xi, yi = np.mgrid[x.min():x.max():200j,
                  y.min():y.max():200j]
zi     = kde(np.vstack([xi.ravel(), yi.ravel()])).reshape(xi.shape)

fig, ax = plt.subplots(figsize=(8, 6))
cf = ax.contourf(xi, yi, zi, levels=15, cmap="Blues")
fig.colorbar(cf, ax=ax, label="Density")
ax.set_title("KDE Density Contour — Dhaka GPS Pings")
ax.set_xlabel("Longitude")
ax.set_ylabel("Latitude")
plt.tight_layout()
plt.savefig("dhaka_kde.png", dpi=150)
\end{lstlisting}

%--------------------------------------------------------------------
\subsection{Challenge 2: Latency and Pre-Aggregation}
\label{subsec:latency}
%--------------------------------------------------------------------

A business analyst using Tableau expects a dashboard filter
(``show Q1 2024, Dhaka region only'') to respond in under one second.
But the underlying fact table may contain 100 billion rows in Hive on HDFS.
A Hive query over the full fact table takes minutes to complete, making
interactive filtering impossible.

Two engineering solutions exist.

\textbf{1. Pre-aggregated data cubes.} A \term{data cube} (also called
an \emph{OLAP cube} or \emph{aggregate table}) materialises the results
of all combinations of GROUP BY queries over a set of dimensions at
ingest time. If the fact table has three dimensions---date,
district, and product category---the data cube pre-computes eight
subtotals: one for each subset of the three dimensions (the
\emph{power set} of dimensions, sometimes called the \emph{cuboid}).

\begin{defbox}{Data Cube}
A \textbf{data cube} over a fact table with $d$ dimensions contains
$2^d$ \emph{cuboids}: one for each subset of dimensions. The base
cuboid aggregates over all dimensions (producing one row); the apex
cuboid contains the full fact table (no aggregation). Materialising
the full cube requires $O(N \cdot 2^d)$ space where $N$ is the fact
table row count, but in practice only a subset of cuboids are
materialised based on query frequency---a process called
\emph{selective materialisation}.

Queries served by the data cube read from a materialised cuboid
rather than the raw fact table, reducing I/O by a factor of
$N / |\text{cuboid}|$---often $10^4$--$10^6 \times$ for typical
analytical workloads.
\end{defbox}

\begin{lstlisting}[style=pyspark, caption={Pre-aggregating a data cube in Spark SQL}]
from pyspark.sql import SparkSession
from pyspark.sql.functions import col, sum as _sum, count, avg

spark = SparkSession.builder.appName("DataCube").getOrCreate()

# Raw fact table: 100 billion rows on HDFS
fact = spark.read.parquet("hdfs:///data/curated/transactions/")

# --- Materialise a data cube with GROUPING SETS ---
# This one Spark job produces all 8 cuboids for 3 dimensions
cube_df = (fact
    .cube("txn_date", "district", "product_category")
    .agg(
        _sum("amount_bdt").alias("total_bdt"),
        count("*").alias("txn_count"),
        avg("amount_bdt").alias("avg_bdt"),
    ))

# Write the cube as Parquet — millions of rows, not billions
(cube_df.write
        .mode("overwrite")
        .partitionBy("txn_date")
        .parquet("hdfs:///data/curated/txn_cube/"))

# Tableau / Superset query: reads from cube (ms), not fact table (min)
district_q1 = spark.read.parquet("hdfs:///data/curated/txn_cube/") \
    .filter(
        (col("txn_date").between("2024-01-01", "2024-03-31")) &
        col("district").isNotNull() &
        col("product_category").isNull()
    ) \
    .orderBy("total_bdt", ascending=False)

district_q1.show(10)
\end{lstlisting}

\textbf{2. Approximate query processing (AQP).} Rather than pre-computing
exact aggregates, \term{approximate query processing} returns a statistically
correct estimate of the aggregate in milliseconds by sampling the table.
\tool{Apache DataSketches} (formerly Yahoo DataSketches) implements a
family of streaming algorithms---HyperLogLog for cardinality, Theta
sketch for set operations, KLL for quantiles---that provide guaranteed
error bounds with sub-linear space. The trade-off: AQP is appropriate
for exploratory dashboards where a 5\% error is acceptable, but not for
financial reconciliation or regulatory reporting where exactness is
mandatory.

%--------------------------------------------------------------------
\subsection{Challenge 3: Streaming Visualization}
\label{subsec:streaming-viz}
%--------------------------------------------------------------------

Many big data applications require visualisations that update
continuously as new events arrive:
\begin{itemize}
  \item An operations dashboard showing the number of active Pathao
        drivers by district, updated every 5 seconds.
  \item A bKash fraud operations centre showing transaction volume
        per minute and active alert count.
  \item A DGHS disease surveillance dashboard showing new case reports
        as hospitals submit their hourly returns.
\end{itemize}

Two delivery mechanisms exist for updating a browser chart with new data.

\begin{defbox}{Poll vs.\ WebSocket Push}
\begin{itemize}
  \item \textbf{Periodic poll (Grafana default)}: The browser makes
        an HTTP request to the backend on a configured interval
        (e.g., every 5 seconds). The backend queries the time-series
        database and returns the latest value. Simplest to implement;
        adds latency of up to one poll interval.
  \item \textbf{WebSocket push}: The server establishes a persistent
        bidirectional connection to the browser. New data points are
        pushed to the browser as they arrive, without waiting for a
        poll. Sub-second update latency; required for real-time financial
        tickers, live election results, and sensor dashboards.
\end{itemize}
\end{defbox}

The standard streaming visualization architecture is:
\begin{enumerate}
  \item \textbf{Kafka}: receives events from producers in real time.
  \item \textbf{Spark Structured Streaming or Flink}: subscribes to
        Kafka topic; applies a tumbling or sliding window aggregation
        (e.g., count of transactions per minute per district).
  \item \textbf{Time-series database} (InfluxDB, Prometheus, TimescaleDB):
        stores the window aggregates; indexed by timestamp.
  \item \textbf{Grafana or Kibana}: polls the TSDB on a configured
        interval (Grafana: minimum 5 seconds) and re-renders the chart.
        For sub-second updates, a custom D3.js frontend with WebSocket
        replaces Grafana.
\end{enumerate}

\begin{figure}[htbp]
\centering
\begin{tikzpicture}[
    box/.style={rectangle, rounded corners=4pt, minimum width=3.0cm,
                minimum height=0.75cm, font=\small\sffamily, align=center},
    arr/.style={-{Stealth[length=4pt,width=4pt]}, thick}
  ]
  \node[box, fill=DEGray!20] (prod) at (0, 6.2)
        {Event Producers\\{\small bKash · Pathao · IoT}};
  \node[box, fill=DEBlue!75, text=white] (kaf) at (0, 5.0)
        {Apache Kafka\\{\small durable event log}};
  \node[box, fill=DEGreen!22] (sp)  at (0, 3.8)
        {Spark Streaming / Flink\\{\small 1-min tumbling window}};
  \node[box, fill=DEAccent!22] (tsdb) at (0, 2.6)
        {Time-Series DB\\{\small InfluxDB / Prometheus}};
  \node[box, fill=DEBlue!20] (graf) at (-2.0, 1.3)
        {Grafana\\{\small poll every 5s}};
  \node[box, fill=DEPurple!18] (d3) at (2.0, 1.3)
        {D3.js + WebSocket\\{\small push, sub-second}};

  \draw[arr, DEGray!50]  (prod)--(kaf);
  \draw[arr, DEBlue!55]  (kaf)--(sp);
  \draw[arr, DEGreen!60] (sp)--(tsdb);
  \draw[arr, DEAccent!60] (tsdb)--(graf);
  \draw[arr, DEAccent!60] (tsdb)--(d3);

  \node[font=\scriptsize\sffamily\itshape, text=DEGray] at (2.8, 2.6)
        {ops dashboards};
  \node[font=\scriptsize\sffamily\itshape, text=DEGray] at (-2.8, 1.3)
        {5s latency};
  \node[font=\scriptsize\sffamily\itshape, text=DEPurple] at (2.8, 1.3)
        {$<$1s latency};
\end{tikzpicture}
\caption{Streaming visualization architecture. Events flow from producers
         through Kafka and a stream processor to a time-series database.
         Grafana polls the TSDB every 5 seconds; a WebSocket-based D3.js
         frontend receives pushed updates for sub-second latency.}
\label{fig:streaming-viz}
\end{figure}

%%====================================================================
\section{The Visualization Tools Landscape}
\label{sec:tools-landscape}
%%====================================================================

The choice of visualization tool determines what is possible, what requires
custom development, and how well the tool integrates with the rest of the
data stack. Six tools dominate the landscape for big data visualisation.

\begin{table}[htbp]
\centering
\renewcommand{\arraystretch}{1.4}
\begin{tabular}{lp{3.0cm}p{2.8cm}p{2.0cm}l}
  \toprule
  \textbf{Tool} & \textbf{Best For}
    & \textbf{Data Sources} & \textbf{Scale} & \textbf{Real-time} \\
  \midrule
  \textbf{Tableau}
    & Business dashboards, drag-and-drop BI
    & SQL, Hive, Spark, CSV, Excel
    & Millions of rows
    & Limited \\
  \textbf{Power BI}
    & Microsoft ecosystem, executive reports
    & Azure, Excel, SQL Server, REST
    & Millions of rows
    & Limited \\
  \textbf{D3.js}
    & Custom interactive web graphics
    & JSON, CSV, REST APIs
    & Flexible (SVG/Canvas)
    & Yes (WebSocket) \\
  \textbf{Grafana}
    & Ops metrics, real-time monitoring
    & InfluxDB, Prometheus, Elasticsearch
    & Millions of data points
    & Yes (5s poll) \\
  \textbf{Apache Superset}
    & Open-source BI on data lakes
    & Presto, Hive, Spark SQL, PG
    & Billions of rows
    & Limited \\
  \textbf{Kibana}
    & Log analytics, search exploration
    & Elasticsearch only
    & Billions of documents
    & Yes \\
  \bottomrule
\end{tabular}
\caption{Visualization tool landscape. Each tool occupies a distinct
         niche; selecting the wrong tool forces expensive workarounds.}
\label{tab:viz-tools}
\end{table}

%--------------------------------------------------------------------
\subsection{Tableau and Power BI: Business Intelligence Platforms}
\label{subsec:tableau-powerbi}
%--------------------------------------------------------------------

\tool{Tableau} (acquired by Salesforce in 2019) is the dominant
business intelligence platform for analyst-facing dashboards. Its drag-and-drop
interface allows non-programmers to build interactive charts, maps, and
cross-tabs from SQL databases, Hive tables, Spark dataframes, and flat files.
Tableau implements Shneiderman's mantra through a drill-down hierarchy
(Year $\to$ Quarter $\to$ Month $\to$ Day), its filter shelf (dynamic
queries), and detail-on-demand tooltips.

Tableau's primary limitation is scale: for tables exceeding
$\sim$50 million rows, Tableau either operates in \emph{extract} mode
(a cached in-memory copy of the data, updated on a schedule) or through
a \emph{live connection} that pushes each query to the underlying database.
For Hive or Presto queries, even a simple filter can take 30--60 seconds
on very large tables, breaking the interactive experience. Pre-aggregated
data cubes in Presto or Spark SQL are the standard mitigation.

\tool{Microsoft Power BI} is the dominant choice in organisations
standardised on the Microsoft Azure ecosystem (Azure Synapse Analytics,
Azure Data Factory, Office 365). Its DAX formula language for calculated
columns and measures is more expressive than Tableau for certain
analytical patterns, but its connector ecosystem is smaller and its
performance at big data scale is similar to Tableau's.

%--------------------------------------------------------------------
\subsection{D3.js: Data-Driven Documents}
\label{subsec:d3}
%--------------------------------------------------------------------

\tool{D3.js} (\textbf{D}ata-\textbf{D}riven \textbf{D}ocuments) is a
JavaScript library for binding data to DOM (Document Object Model)
elements in a web browser and applying transformations based on data
values. Created by Mike Bostock, Vadim Ogievetsky, and Jeffrey Heer at
Stanford and released in 2011, D3 is the foundation of the New York Times
interactive graphics, The Guardian's data journalism, and the Reuters
COVID-19 tracker.

D3's design philosophy is minimal abstraction: it provides general
geometric and visual primitives (SVG shapes, scales, axes, transitions)
but does not prescribe chart types. Every element of a D3 chart is
created by writing JavaScript that explicitly binds a data array to
SVG elements:

\begin{lstlisting}[style=pyspark, caption={D3.js: Tufte-style bar chart (data-driven SVG)}]
/* Tufte-compliant D3 bar chart — data-ink ratio maximised
   No gridlines, no background, direct data labels, 0-based y-axis */

const data = [
  { district: "Dhaka",       cases: 4217 },
  { district: "Chittagong",  cases: 1803 },
  { district: "Khulna",      cases:  892 },
  { district: "Sylhet",      cases:  671 },
  { district: "Rajshahi",    cases:  544 },
];

const margin = { top: 20, right: 60, bottom: 40, left: 90 };
const W = 560 - margin.left - margin.right;
const H = 320 - margin.top  - margin.bottom;

const svg = d3.select("#chart")
  .append("svg")
    .attr("width",  W + margin.left + margin.right)
    .attr("height", H + margin.top  + margin.bottom)
  .append("g")
    .attr("transform", `translate(${margin.left},${margin.top})`);

const x = d3.scaleLinear()
  .domain([0, d3.max(data, d => d.cases)])
  .nice()
  .range([0, W]);

const y = d3.scaleBand()
  .domain(data.map(d => d.district))
  .range([0, H])
  .padding(0.35);

// Bars — single flat colour (no gradient, no shadow)
svg.selectAll(".bar")
   .data(data)
   .join("rect")
     .attr("class", "bar")
     .attr("x",       0)
     .attr("y",       d => y(d.district))
     .attr("width",   d => x(d.cases))
     .attr("height",  y.bandwidth())
     .attr("fill",    "#003c78");   // DEBlue

// Direct data labels (no legend needed)
svg.selectAll(".label")
   .data(data)
   .join("text")
     .attr("x",           d => x(d.cases) + 4)
     .attr("y",           d => y(d.district) + y.bandwidth() / 2 + 4)
     .text(d => d.cases.toLocaleString())
     .attr("font-size",   11)
     .attr("fill",        "#505050");

// Minimal y-axis: district names only
svg.append("g").call(d3.axisLeft(y).tickSize(0));

// Minimal x-axis: no domain line, light ticks only
svg.append("g")
   .attr("transform", `translate(0,${H})`)
   .call(d3.axisBottom(x).ticks(5).tickSizeOuter(0));

/* Tooltip (Shneiderman: details-on-demand) */
const tooltip = d3.select("body").append("div")
  .style("position", "absolute")
  .style("background", "#fff")
  .style("border", "1px solid #ccc")
  .style("padding", "6px 10px")
  .style("font-size", "12px")
  .style("pointer-events", "none")
  .style("opacity", 0);

svg.selectAll(".bar")
   .on("mouseover", (event, d) => {
     tooltip.style("opacity", 1)
            .html(`<b>${d.district}</b><br>Cases: ${d.cases.toLocaleString()}`);
   })
   .on("mousemove", event => {
     tooltip.style("left",  (event.pageX + 10) + "px")
            .style("top",   (event.pageY - 28) + "px");
   })
   .on("mouseout", () => tooltip.style("opacity", 0));
\end{lstlisting}

%--------------------------------------------------------------------
\subsection{Grafana and Apache Superset: Open-Source Platforms}
\label{subsec:grafana-superset}
%--------------------------------------------------------------------

\tool{Grafana} is the dominant open-source tool for operations and
infrastructure monitoring dashboards. It natively integrates with
\tool{InfluxDB} (time-series database optimised for metrics),
\tool{Prometheus} (pull-based monitoring used in Kubernetes deployments),
and \tool{Elasticsearch} (for log aggregates). Grafana implements
Shneiderman's zoom-and-filter via its time-range selector (click-drag
on any chart to zoom into a sub-interval) and variable dropdowns
(dynamic query parameters). Its stat panel widget renders a number with
a sparkline---a direct implementation of Tufte's sparkline concept.

\tool{Apache Superset} (open-source, Apache Software Foundation) is
the standard open-source BI tool for organisations that run their data
platform on Presto, Trino, Apache Hive, or Spark SQL. Its SQL Lab
feature allows analysts to write arbitrary SQL and render the results
as charts; its dashboard builder allows interactive filter boxes and
cross-chart drill-down. For organisations that cannot afford or prefer
not to use Tableau, Superset provides comparable BI functionality with
direct connectivity to the Hadoop/Spark ecosystem.

%%====================================================================
\section{Research Spotlight}
\label{sec:ch9-research}
%%====================================================================

\begin{researchspotlight}{Tufte (1983) · Shneiderman (1996) · Bostock,
Ogievetsky \& Heer (2011) --- The Three Pillars of Data Visualization}

\textbf{Three works that together define the field:}

\medskip
\textbf{Work 1: Tufte, E.\,R. (1983). \textit{The Visual Display of
Quantitative Information.} Graphics Press.} ($>$20{,}000 citations.)

Tufte analysed 4{,}000 charts from science, business, and journalism
publications and derived the data-ink ratio principle: of all the ink
used in a graphic, maximise the fraction that encodes data. He coined
``chartjunk'' to describe ink that encodes no data (grids, 3-D effects,
gradients, redundant legends) and documented the lie factor---a metric
for quantifying graphic deception from axis truncation or scale distortion.
His concept of small multiples described a grid of same-scale same-axes
charts as the most efficient form of comparative display, and his sparklines
(formalised in his 2006 follow-up) described word-sized trend indicators
now ubiquitous in Grafana, Bloomberg terminals, and spreadsheets. Tufte's
book is unique in data engineering in that it both states the principle and
implements it beautifully in the book's own typography and layout: the book
itself has a high data-ink ratio. Every modern dashboard framework---from
Tableau's clean minimal default theme to D3's blank-canvas philosophy---is
a direct application of Tufte's principles.

\medskip
\textbf{Work 2: Shneiderman, B. (1996). ``The Eyes Have It: A Task by
Data Type Taxonomy for Information Visualizations.'' \textit{Proc.\ IEEE
Visual Languages}, pp.~336--343.} ($>$5{,}000 citations.)

Shneiderman identified seven types of data (linear, 2-D planar,
3-D volumetric, temporal, multi-dimensional, tree, network) and seven
tasks that users perform on visualisations (overview, zoom, filter,
details-on-demand, relate, history, extract). He proposed the
``visual information seeking mantra''---overview first, zoom and filter,
then details on demand---as a universal three-step design pattern for
interactive information visualisation. He also invented the \emph{treemap}
(1991) and the \emph{dynamic query} concept (sliders and checkboxes that
update the visualisation in $<$100~ms). The 1996 paper's taxonomy
directly shapes how modern BI platforms are designed: Tableau's hierarchy
drill-down implements ``overview $\to$ zoom''; Grafana's variable dropdowns
implement ``filter''; hover tooltips implement ``details-on-demand.''
The mantra is the single most widely implemented design pattern in
data visualisation software.

\medskip
\textbf{Work 3: Bostock, M., Ogievetsky, V., \& Heer, J. (2011). ``D$^3$:
Data-Driven Documents.'' \textit{IEEE Transactions on Visualization and
Computer Graphics}, 17(12), 2301--2309.} ($>$3{,}000 citations.)

Bostock, Ogievetsky, and Heer introduced D3.js---a JavaScript library
that binds data arrays to DOM elements (SVG shapes, HTML elements) and
applies visual transformations based on data values. The key innovation
over previous web charting libraries (Protovis, Raphaël) was \emph{general
selections}: any set of DOM elements can be bound to any array; the
enter-update-exit pattern handles data-driven creation, modification, and
removal of elements. D3 does not prescribe chart types; instead it
provides general primitives (scales, axes, path generators, force
simulations, geographic projections) from which any chart can be assembled.
This philosophy---maximum expressive power at the cost of requiring more
code---directly reflects Tufte's principle that the graphic designer must
control every pixel of ink. D3 became the production tool of the New York
Times Graphics Desk, The Guardian, Reuters, the FT's Visual Journalism
team, and the US 2020 Census Bureau's interactive data explorer. It is the
most widely cited visualisation systems paper of the 2010s and is used in
the graduate visualisation curriculum at Harvard (CS171), MIT (6.859),
and Stanford (CS448B).
\end{researchspotlight}

%%====================================================================
\section{Case Study: The New York Times Graphics Desk}
\label{sec:nyt-case}
%%====================================================================

\begin{casestudy}{New York Times Graphics Desk --- Interactive Data
Journalism at 6 Million Readers per Day}

\textbf{Context.} The New York Times Graphics Desk is a team of
approximately 50 journalists, designers, and engineers who produce
interactive data visualisations for nytimes.com---reaching 6~million
readers per day. The desk is responsible for some of the most widely
read and most cited data visualisations of the past decade, including
the 2012 US presidential election maps, the 2016 election needle, and
the COVID-19 tracker that was loaded 800~million times in 2020.

Mike Bostock---co-creator of D3.js---worked at the NYT Graphics Desk
from 2010 to 2015, during which time he created D3.js to replace the
desk's previous tool (Protovis). D3.js was born at the NYT not as a
research project but as a production engineering need: the desk required
pixel-level control over chart layout, custom geographic projections,
and animated transitions between data states that no existing library
provided.

\medskip
\textbf{Technical stack.}
\begin{center}
\small
\renewcommand{\arraystretch}{1.25}
\begin{tabular}{ll}
  \toprule
  \textbf{Component} & \textbf{Technology} \\
  \midrule
  Data collection and cleaning  & Python (pandas, requests, BeautifulSoup) \\
  Geographic processing         & PostGIS, GDAL, Mapbox \\
  Statistical analysis          & R (ggplot2, tidyverse) \\
  Interactive web charts        & D3.js (SVG and Canvas) \\
  Static print graphics         & Adobe Illustrator + ai2html \\
  Collaborative prototyping     & Observable Notebooks \\
  Data publication              & GitHub CSV repository \\
  Content delivery              & nytimes.com CDN (AWS CloudFront) \\
  \bottomrule
\end{tabular}
\end{center}

\medskip
\textbf{COVID-19 Tracker (2020--2022): a big data visualisation at journalism scale.}

From March 2020 the desk built and maintained the most-read COVID-19
data tracker in the English-speaking world. The data engineering
challenges were representative of any big data pipeline at scale:

\begin{enumerate}
  \item \textbf{Data acquisition at scale.} The tracker aggregated case
        and death data from the US Centers for Disease Control (CDC),
        50 state health departments, and 3{,}000 county health departments
        ---each with different update schedules, different column names,
        and different definitions of ``confirmed case.'' Python scrapers
        ran hourly; a pandas pipeline normalised schemas and detected
        anomalies (negative daily case counts from retroactive
        corrections; impossible date values; duplicate state-county records).

  \item \textbf{The rolling average as a design decision.} Raw daily
        counts exhibited strong day-of-week artefacts: fewer cases were
        reported on weekends because fewer testing sites operated.
        The desk chose a 7-day rolling average to smooth these artefacts,
        encoding the average as the primary line and the raw counts as a
        lighter bar chart behind it. This is a \emph{journalistic} choice
        with statistical implications: the 7-day window delays the
        detection of a surge by 3--4 days compared to the raw count.

  \item \textbf{Tufte principles in production.} The 50-state
        small-multiples chart (all states on a single page, same scale,
        same axes) allowed readers to compare New York, Texas, and
        California in one glance without animation. No gridlines appeared
        on the case charts; only a single light horizontal reference line
        at the prior peak provided context. Hover tooltips provided
        exact date, raw count, and 7-day average (Shneiderman's
        details-on-demand). The colour scale was tested for
        colour-blindness (deuteranopia and protanopia) using the
        \texttt{chroma.js} simulator.

  \item \textbf{Accessibility at 800 million page loads.} Every chart
        had a data-table alternative for screen readers. SVG title
        and desc elements described each chart's content to assistive
        technology. This was not optional: Section 508 of the US
        Rehabilitation Act requires federal government and major news
        organisations to make digital content accessible.

  \item \textbf{GitHub as data infrastructure.} The NYT published its
        cleaned CSV data daily on GitHub (\texttt{nytimes/covid-19-data}).
        Within 30 days of the tracker's launch, over 1{,}000 derivative
        dashboards had been built by universities, local governments,
        and independent developers worldwide---all using the NYT's
        cleaned dataset. This transformed the NYT from a single publisher
        into a data infrastructure provider.
\end{enumerate}

\medskip
\textbf{Engineering lessons.}
\begin{enumerate}
  \item \textbf{Data quality is the invisible half of visualisation.}
        The hardest engineering work on the COVID tracker was schema
        normalisation across 50 health departments---not SVG rendering.
        A visualisation pipeline without a robust data quality gate
        produces misleading charts with high production values.
  \item \textbf{Small multiples beat animation for comparative analysis.}
        The 50-state grid was more actionable than an animated spreading-map
        because it allowed simultaneous comparison of all states
        without requiring readers to hold past frames in working memory.
  \item \textbf{Accessibility is an engineering requirement, not a
        nicety.} At 6~million readers per day, even 1\% of readers
        with visual impairments represents 60{,}000 people per day who
        would be excluded by inaccessible charts.
  \item \textbf{Rolling averages are editorial decisions.}
        The choice of window size changes the story the chart tells.
        A 3-day window reveals surges earlier but amplifies noise;
        a 14-day window is more stable but delays detection. These
        decisions have public health consequences and must be documented
        transparently in the chart's methodology notes.
\end{enumerate}
\end{casestudy}

%%====================================================================
\section*{Chapter Summary}
\label{sec:ch9-summary}
%%====================================================================

Effective data visualization is not a presentation afterthought: it is
the first quality check on any analytical result and the primary interface
between data infrastructure and human decision-making.

Tufte's data-ink ratio provides an actionable metric for graphical
efficiency: maximise the fraction of ink that encodes data; erase
everything else. His lie factor quantifies graphical deception from
axis truncation. Small multiples and sparklines solve the
comparative-display problem without animation. Shneiderman's mantra
(overview $\to$ zoom/filter $\to$ details-on-demand) is the universal
interaction design pattern for interactive information visualisation,
implemented in every major dashboard platform.

Big data introduces three engineering challenges not present with small
datasets: overplotting (solved by hexbin aggregation and KDE contours),
latency (solved by pre-aggregated data cubes and approximate query
processing), and streaming (solved by the Kafka $\to$ stream processor
$\to$ TSDB $\to$ Grafana/WebSocket architecture). Tool selection---Tableau
for analyst-facing BI, D3.js for custom interactive journalism and
production web graphics, Grafana for real-time operational monitoring,
Apache Superset for open-source lake-native analytics---determines what
is achievable without custom development.

The New York Times COVID-19 tracker demonstrated that at journalistic
scale the data quality pipeline is harder than the rendering pipeline,
and that responsible visualization design---rolling-average window
selection, colour-blindness testing, accessibility requirements---carries
ethical weight proportional to audience size.

%%====================================================================
\section*{Exercises}
\label{sec:ch9-exercises}
%%====================================================================

\exercisesection{Short Questions}

\sq What is Anscombe's Quartet? What property of summary statistics does
it illustrate, and what does this imply about the role of visualization
in data analysis?

\sq Define Tufte's data-ink ratio. Identify three specific chart elements
that commonly lower the data-ink ratio and should be removed.

\sq A bar chart shows four annual revenue figures. The data changed from
\$82M in 2022 to \$90M in 2023---a 9.8\% increase. The chart's $y$-axis
starts at \$80M. The bar for 2023 appears five times taller than the
bar for 2022. Calculate Tufte's lie factor.

\sq What is a ``small multiples'' chart? Give an example of a small
multiples grid that would be more effective than an animated chart for
communicating the same information.

\sq State Shneiderman's visual information seeking mantra in full.
Give one example of how each of the three steps is implemented in Tableau,
in Grafana, and in a D3.js chart.

\sq What is overplotting? At what approximate number of data points does
it typically become a problem in a standard scatterplot?

\sq Describe the hexbin technique for addressing overplotting. What does
it show that alpha transparency cannot?

\sq What is a data cube? If a fact table has four dimensions---date,
district, category, and channel---how many cuboids does the full data
cube contain?

\sq Explain the difference between a \emph{periodic poll} and a
\emph{WebSocket push} architecture for streaming visualization. Give
a use case where each is appropriate.

\sq Which visualization tool is most appropriate for: (a) a real-time
operations dashboard monitoring Kafka consumer lag; (b) a business analyst
building a quarterly sales report in a company that uses Azure; (c) a data
journalist building a custom interactive choropleth map for publication
on a news website?

\sq What is the Shneiderman treemap? Explain how area and colour are
used to encode two variables simultaneously.

\sq Why does Tufte argue that 3-D effects on bar charts lower the
data-ink ratio \emph{and} increase the lie factor?

\exercisesection{Analytical Problems}

\ap \textbf{Data-Ink Ratio Audit.}
A data engineer builds a Grafana dashboard for a Kafka pipeline showing
hourly message throughput per topic. The current chart has the following
elements: a 3-D bar chart; a $y$-axis that starts at 500{,}000 messages/hour
(when the true minimum is 480{,}000); a thick black border around the chart
area; a gradient fill on each bar (dark blue at bottom, light blue at top);
a legend listing all 8 topic names in a box overlapping the chart; and
heavy horizontal gridlines every 10{,}000 messages.
\begin{enumerate}[(a)]
  \item Identify every element that violates Tufte's data-ink principle,
        classify each as chartjunk or a lie-factor violation, and explain
        what information is lost by removing it (if any).
  \item Compute the approximate lie factor for the truncated $y$-axis
        if the true range is 480{,}000--560{,}000 messages/hour but the
        chart displays 500{,}000--560{,}000. The highest bar appears to be
        $5\times$ the height of the lowest bar on screen.
  \item Redesign the chart: describe each element of the corrected
        version, applying Tufte's principles.
\end{enumerate}

\ap \textbf{Shneiderman Dashboard Design.}
A Pathao data scientist needs to build an operations dashboard for the
Dhaka fleet management team showing:
\begin{itemize}
  \item City-wide active driver count by district (64 districts)
  \item Trip completion rate per hour over the past 7 days
  \item Distribution of trip duration (in minutes) for completed trips
  \item Geographic density of trip pickups (GPS coordinates)
\end{itemize}
For each requirement:
(a) identify Shneiderman's data type;
(b) choose the appropriate chart type, citing the data-type taxonomy;
(c) specify how overview, zoom/filter, and details-on-demand are implemented;
(d) identify which big data challenge (overplotting, latency, or streaming)
applies and how to address it.

\ap \textbf{Streaming Visualization Latency Budget.}
A bKash fraud operations dashboard must update its fraud alert rate
chart within 10 seconds of an event arriving at the source system.
The pipeline is: Mobile app $\to$ Kafka (1 partition, 50~MB/s)
$\to$ Spark Structured Streaming (1-minute tumbling window)
$\to$ InfluxDB $\to$ Grafana (WebSocket push).

Allocate the 10-second latency budget across the pipeline components.
For each component: (a) state the typical latency contribution; (b) identify
one configuration parameter that can reduce it; (c) identify the trade-off
of reducing it.

\ap \textbf{Data Cube Design.}
A national electricity utility has a 5-year fact table with 120 billion
rows and four dimensions: \texttt{meter\_id} (10 million distinct values),
\texttt{hour\_ts} (43{,}800 hours over 5 years), \texttt{district}
(64 districts), and \texttt{tariff\_band} (5 bands).
\begin{enumerate}[(a)]
  \item How many cuboids are in the full data cube? How many are practically
        useful for a dashboard (explain which dimensions are too high
        cardinality to materialise)?
  \item Design the materialised cuboids for a dashboard that must serve:
        (i) hourly consumption per district for any 30-day window;
        (ii) monthly total consumption per tariff band; (iii) peak-hour
        demand by district.
  \item Estimate the storage size (in GB) of the hourly-by-district
        cuboid assuming 30 bytes per row.
\end{enumerate}

\exercisesection{Design Problems}

\dpr \textbf{DGHS Disease Surveillance Dashboard.}
The Directorate General of Health Services (DGHS), Bangladesh, needs a
public health dashboard communicating:
\begin{enumerate}[(a)]
  \item National COVID-19 confirmed case counts by district over 24 months,
        using a 7-day rolling average.
  \item Real-time hospitalisation rate updated every 5 minutes from hospital
        daily reporting forms submitted via a REST API.
  \item Demographic breakdown (age group $\times$ gender) of confirmed cases
        as a proportion of population.
  \item A ``situation map'' showing the 7-day incidence rate per 100{,}000
        population for all 64 districts.
\end{enumerate}
Design the complete dashboard: specify the chart type and justification
for each panel; describe how Shneiderman's mantra is implemented for each
panel; describe the data pipeline (ingestion, aggregation, storage, and
rendering tools); identify and address the big data visualization challenge
for each panel; and describe two accessibility requirements.

\dpr \textbf{Streaming bKash Operations Dashboard.}
Design a real-time operations dashboard for bKash (10 million transactions/day,
peak 500 transactions/second) that shows:
\begin{enumerate}[(a)]
  \item Total transaction volume (BDT) per minute, last 60 minutes,
        updated every 5 seconds.
  \item Number of active fraud alerts, updated every 30 seconds.
  \item Geographic heatmap of transaction density across Dhaka (GPS
        coordinates from merchant terminal locations).
  \item Top-5 error codes by frequency in the past 15 minutes.
\end{enumerate}
Specify: the Kafka topic structure and partition count for each stream;
the Spark Structured Streaming window and aggregation for each metric;
the TSDB schema for each metric; the Grafana panel type and refresh interval;
and any pre-aggregation required to achieve the specified latency.

\exercisesection{Programming Exercises}

\pe \textbf{Overplotting Solutions in Python.}
Given a simulated dataset of 2 million Pathao GPS trip origin points
(longitude, latitude) distributed across five Dhaka zones:
\begin{enumerate}[(a)]
  \item Plot the raw scatterplot with $\alpha = 0.005$. Describe the
        overplotting effect observed.
  \item Create a hexbin plot with \texttt{gridsize=100} and a sequential
        blue colour scale. Describe what information is now visible.
  \item Compute and plot a 2-D KDE contour using \texttt{scipy.stats.gaussian\_kde}
        on a 10{,}000-point random subsample. Identify the location of the
        three highest-density zones.
  \item Create a 2$\times$2 panel figure (overview in panel 1; one Dhaka zone
        zoomed in panels 2--4 with hexbin). Label each panel with its
        Shneiderman interaction level (overview, zoom, detail).
\end{enumerate}

\begin{lstlisting}[style=pyspark, caption={PE9.1: GPS overplotting simulation}]
import numpy as np
import matplotlib
matplotlib.use("Agg")
import matplotlib.pyplot as plt
from scipy.stats import gaussian_kde

rng = np.random.default_rng(2024)

# Five Dhaka zones: (lon_center, lat_center, sigma, n_points)
zones = [
    (90.400, 23.810, 0.020, 600_000),  # Gulshan / Banani
    (90.380, 23.740, 0.015, 500_000),  # Dhanmondi
    (90.350, 23.790, 0.018, 400_000),  # Mirpur
    (90.420, 23.870, 0.012, 300_000),  # Uttara
    (90.410, 23.720, 0.010, 200_000),  # Motijheel
]

lons, lats = [], []
for lon_c, lat_c, sigma, n in zones:
    lons.append(rng.normal(lon_c, sigma, n))
    lats.append(rng.normal(lat_c, sigma, n))

lons = np.concatenate(lons)
lats = np.concatenate(lats)

fig, axes = plt.subplots(2, 2, figsize=(14, 12))

# Panel 1 — Overview: raw scatter (overplotted)
axes[0,0].scatter(lons, lats, s=0.003, alpha=0.004, c="steelblue")
axes[0,0].set_title("(a) Scatterplot — overplotted (Shneiderman: overview)")

# Panel 2 — Zoom: hexbin full city
hb = axes[0,1].hexbin(lons, lats, gridsize=100, cmap="Blues", mincnt=1)
fig.colorbar(hb, ax=axes[0,1], label="Trip count")
axes[0,1].set_title("(b) Hexbin — density readable (Shneiderman: zoom)")

# Panel 3 — Zoom: KDE contour on subsample
idx     = rng.choice(len(lons), 10_000, replace=False)
kde     = gaussian_kde(np.vstack([lons[idx], lats[idx]]), bw_method=0.08)
xi, yi  = np.mgrid[lons.min():lons.max():150j, lats.min():lats.max():150j]
zi      = kde(np.vstack([xi.ravel(), yi.ravel()])).reshape(xi.shape)
cf = axes[1,0].contourf(xi, yi, zi, levels=12, cmap="Blues")
fig.colorbar(cf, ax=axes[1,0], label="Density")
axes[1,0].set_title("(c) KDE contour — shape revealed (Shneiderman: zoom)")

# Panel 4 — Detail: Gulshan zone only
mask = (lons > 90.36) & (lons < 90.44) & (lats > 23.79) & (lats < 23.83)
axes[1,1].hexbin(lons[mask], lats[mask], gridsize=60, cmap="Reds", mincnt=1)
axes[1,1].set_title("(d) Gulshan zone detail (Shneiderman: detail)")

for ax in axes.ravel():
    ax.set_xlabel("Longitude")
    ax.set_ylabel("Latitude")

plt.tight_layout()
plt.savefig("dhaka_overplotting.png", dpi=150, bbox_inches="tight")
print("Saved: dhaka_overplotting.png")
\end{lstlisting}

\pe \textbf{Pre-Aggregated Data Cube (PySpark).}
Using PySpark, build a pre-aggregated data cube for a simulated 100
million-row e-commerce fact table with columns: \texttt{order\_date},
\texttt{district}, \texttt{category}, and \texttt{revenue\_bdt}.
\begin{enumerate}[(a)]
  \item Generate 100 million synthetic rows using Spark's built-in
        random functions.
  \item Use \texttt{DataFrame.cube()} to materialise all cuboids for the
        three dimensions \texttt{(order\_date, district, category)}.
  \item Write the cube as Parquet, partitioned by \texttt{order\_date}.
  \item Simulate three Tableau/Superset queries (filter by date range,
        filter by district, group by category) against the cube and
        measure their execution time using \texttt{time.perf\_counter()}.
        Compare to the same queries against the raw fact table.
  \item Report the size of the materialised cube in MB and the speedup
        factor for each query.
\end{enumerate}

\pe \textbf{Streaming Dashboard Simulation (Python + InfluxDB).}
Build a simulated streaming pipeline using Python threads:
\begin{enumerate}[(a)]
  \item Thread 1 (``producer''): generates 100 synthetic bKash
        transaction events per second, writing them to a Python queue.
  \item Thread 2 (``stream processor''): reads from the queue; aggregates
        over a 10-second tumbling window; computes total BDT, transaction
        count, and fraud flag count per window.
  \item Thread 3 (``TSDB writer''): writes each window's aggregates
        to an InfluxDB local instance (using the \texttt{influxdb-client}
        Python library) with a \texttt{bkash\_metrics} measurement and
        fields \texttt{total\_bdt}, \texttt{txn\_count},
        \texttt{fraud\_count}.
  \item Configure a local Grafana instance to display the three metrics
        as a time-series chart with a 10-second auto-refresh interval.
  \item Run the pipeline for 5 minutes and capture a screenshot
        showing the live chart updating.
\end{enumerate}

\begin{lstlisting}[style=pyspark, caption={PE9.3: Streaming aggregation and InfluxDB write}]
import queue, time, random, threading
from datetime import datetime, timezone
from influxdb_client import InfluxDBClient, Point
from influxdb_client.client.write_api import SYNCHRONOUS

INFLUX_URL    = "http://localhost:8086"
INFLUX_TOKEN  = "mytoken"
INFLUX_ORG    = "cse4345"
INFLUX_BUCKET = "bkash_metrics"
WINDOW_SEC    = 10

event_queue: queue.Queue = queue.Queue(maxsize=5000)

def producer():
    """Generates 100 synthetic events/sec."""
    while True:
        for _ in range(100):
            event_queue.put({
                "amount":    round(random.uniform(50, 5000), 2),
                "is_fraud":  random.random() < 0.003,   # 0.3% fraud rate
                "ts":        time.time(),
            })
        time.sleep(1)

def stream_processor(write_api):
    """Aggregates events in 10-second tumbling windows."""
    window_start = time.time()
    total_bdt, txn_count, fraud_count = 0.0, 0, 0

    while True:
        try:
            event = event_queue.get(timeout=0.05)
            total_bdt   += event["amount"]
            txn_count   += 1
            fraud_count += int(event["is_fraud"])
        except queue.Empty:
            pass

        now = time.time()
        if now - window_start >= WINDOW_SEC:
            # Write window aggregate to InfluxDB
            point = (Point("bkash_metrics")
                     .field("total_bdt",   total_bdt)
                     .field("txn_count",   txn_count)
                     .field("fraud_count", fraud_count)
                     .time(datetime.now(timezone.utc)))
            write_api.write(bucket=INFLUX_BUCKET, org=INFLUX_ORG, record=point)
            print(f"Window: {txn_count} txns | "
                  f"BDT {total_bdt:,.0f} | "
                  f"Fraud {fraud_count}")
            total_bdt, txn_count, fraud_count = 0.0, 0, 0
            window_start = now

if __name__ == "__main__":
    client    = InfluxDBClient(url=INFLUX_URL, token=INFLUX_TOKEN, org=INFLUX_ORG)
    write_api = client.write_api(write_options=SYNCHRONOUS)

    t_prod = threading.Thread(target=producer, daemon=True)
    t_proc = threading.Thread(target=stream_processor,
                              args=(write_api,), daemon=True)
    t_prod.start()
    t_proc.start()

    print("Pipeline running. Press Ctrl+C to stop.")
    time.sleep(300)   # run for 5 minutes
\end{lstlisting}

%%====================================================================
\section*{Further Reading}
\label{sec:ch9-further}
%%====================================================================

\begin{itemize}
  \item \textbf{\cite{Tufte1983}} --- Self-published masterpiece that
        founded the analytical visualization discipline. Chapters 1--4
        cover data-ink ratio, chartjunk, and the lie factor with
        hundreds of printed examples.
  \item \textbf{\cite{Shneiderman1996}} --- The original mantra paper.
        Five pages. Read it in full; every page has a principle directly
        applicable to dashboard design.
  \item \textbf{\cite{Bostock2011}} --- The D3 paper. Section 3 (Design
        Considerations) explains why D3 provides general primitives rather
        than chart types, and how the enter-update-exit pattern works.
        Reading this alongside the D3 API documentation is the fastest
        path to production D3 fluency.
  \item \textbf{\cite{Acharya2015}} --- Chapter 10 covers data
        visualisation tools for big data, including Tableau, R, and Python
        with a practical emphasis suited to the course level.
  \item \textbf{Wilke, C.\,O. (2019). \textit{Fundamentals of Data
        Visualization.} O'Reilly.} A modern, freely available book that
        extends Tufte's principles with colour theory, perceptual accuracy
        research, and R ggplot2 implementation examples.
        (Available at \texttt{clauswilke.com/dataviz}.)
  \item \textbf{Bostock, M. (2022). Observable Plot documentation.}
        Bostock's successor to D3 for exploratory data analysis in the
        browser; simpler API while preserving the data-binding philosophy.
\end{itemize}
